\subsection{自然整数集}

定义1. 如果一个集合不是有限的，则称其为无限的。

特别地，一个基数如果不是整数，则它是无限的。

关系“存在一个无限集”蕴含关系“x是整数”是收集性的（第二章，§1，第4号）；因为如果a是一个无限基数且n是任意整数，我们不可能有 \(a \leqslant n\) （§4，第2号，命题2）。因此我们有 n \textless{} a 对所有整数 n 均成立，这表明整数集合 \textless{} a（§ 3，第 2 号，定理 1 后的注记）包含所有整数。反之，若关系“x 是整数”是收集性的，则整数集合 E 是一个无限集。对每个整数 n，区间 {[}0, n{]} 是E的具有 \(n + 1\) 个元素的子集（§ 5，第 3 号，命题 5）。因此 Card (E) \(\geqslant n + 1 > n\) 。但要说 \(\text{Card}(E) \neq n\) 对每个整数 n 成立，就意味着 E 是无限的。

我们现在引入以下公理：

A5（“无穷公理”。）存在一个无限集。

目前尚不知道此公理是否可以从先前引入的公理和公理模式中推导出来；尽管该问题尚未最终解决，但可以推测此公理与其他公理是独立的。

前面的论述于是证明了以下定理：

定理 1. “x 是整数”这一关系是收集性的。

我们用 N 表示整数集（在必要时为避免歧义，也称为“自然整数集”）。N 的基数记为 \(\aleph_0\) 。每当 \(\mathbf{N}\) 被视为有序集时，所考虑的总是§3第2号中定义的序（称为通常序），除非另有明确说明。

定义2. 一个序列（相应地，集合E的元素序列）是一个族（相应地，E的元素族），其指标集是N的一个子集。如果其指标集是N的一个无限子集，则该序列称为无限序列。

让 \(P\{n\}\) 为一个关系，并设I表示使得 \(P\{n\}\) 为真的整数n的集合。则I是N的一个子集。一个序列 \((x_n)_{n \in \mathbf{I}}\) 有时被写成 \((x_n)_{P\{n\}}\) ，以及 \(x_n\) 称为该序列的第n项。指标集为整数集 \(n \geqslant k\) 的序列常被写成 \((x_n)_{k \leqslant n}\) 或 \((x_n)_{n \geqslant k}\) ，或者如果k=0或k=1，甚至只写成 \((x_n)\) 。在相同的

条件下，例如，记号 \(\prod_{\mathbf{P}\nmid n}\mathbf{X}_n\) 并且 \(\prod_{n = k}\mathbf{X}_n\) 用于表示集合序列 \((\mathbf{X}_n)_{n\in \mathbf{I}}\) 的乘积，并且对于并、交、基数乘积和基数和也有类似的记号。

序列的每个子族都是一个序列，称为给定序列的子序列。

两个序列 \((x_{n})_{n\in I}, (y_{n})_{n\in I}\) 若具有相同的指标集，则当存在一个置换 \(f\) （指标集 \(I\) 的置换），使得 \(x_{f(n)} = y_{n}\) 对于所有 \(n\in I\) 时，称它们仅在其项的次序上有所不同.

多重序列是一个族，其指标集是乘积 \(N^{p}\) （p为整数）的一个子集（p=2时称为“二重序列”，p=3时称为“三重序列”，依此类推）。

设 I 是一个与 N 等势的集合，并设 f 是 N 到 I 的一个双射。对于每个由集合 I 索引的族 \((x_{t})_{t\in I}\) ，序列 \(n\to x_{f(n)}\) 被称为按 f 所定义的顺序排列该族 \((x_{t})_{t\in I}\) 而得到的序列。以这种方式对应于 N 到 I 的两个不同双射的序列，仅在其项的次序上有所不同。对于由具有 n 个元素的集合 I 索引的有限族，我们同样可以定义一个有限序列，以 \([1,n]\) 或 \([0,n-1]\) 作为指标集，通过按照由这两个区间之一到I的双射所定义的顺序排列该族。

